\documentclass[a4j,twocolumn,11pt]{jarticle}
\usepackage{view2025}
\usepackage[hidelinks,dvipdfmx]{hyperref}

\makeatletter
%%%タイトル調整
\setcounter{totalnumber}{3}
\def\id#1{\def\@id{#1}}
\def\@maketitle{%
\vspace{-4.3em}
\begin{center}%
{\@title \par}% タイトル
{\@author}% 著者
\end{center}%
%\par\vskip 1.5em
}
%%%参考文献調整
\renewenvironment{thebibliography}[1]
{\section*{\refname\@mkboth{\refname}{\refname}}%
  \list{\@biblabel{\@arabic\c@enumiv}}%
       {\settowidth\labelwidth{\@biblabel{#1}}%
       \leftmargin\labelwidth
       \advance\leftmargin\labelsep
       \setlength\itemsep{-0.5zh}
       \setlength\baselineskip{9pt}
       \@openbib@code
       \usecounter{enumiv}%
       \let\p@enumiv\@empty
       \renewcommand\theenumiv{\@arabic\c@enumiv}}%
  \sloppy
  \clubpenalty4000
  \@clubpenalty\clubpenalty
  \widowpenalty4000%
  \sfcode`\.\@m}
 {\def\@noitemerr
   {\@latex@warning{Empty `thebibliography' environment}}%
  \endlist}
\makeatother

\title{
\vspace{60pt}
\LARGE{
\textbf{キャラクターの基準印象と服装変換要因に基づく\\
ギャップ萌えの感性分析と予測}}
}

\author{
\vspace{1.5em}
○杉原 明音 $\dagger$， 山口 熙己 $\dagger$，植木 一也 $\dagger$\\
\vspace{0.5em}
$\dagger$：明星大学 情報学部 情報学科\\
\textrm{23j5062@stu.meisei-u.ac.jp, 23j5126@stu.meisei-u.ac.jp, kazuya.ueki@meisei-u.ac.jp} \\
}
\date{} %日付を表示しない
\pagestyle{empty}

\setstretch{1.1} %%%行間隔調整用：適宜調整してください．
%\pagestyle{empty}
\begin{document}
%---タイトル---
\twocolumn[
\maketitle
\thispagestyle{empty} %maketitleをいじっている関係でここでページ番号を消去
\vspace{-2em}
\begin{abstract}
 \noindent{概要：}生成AIによりキャラクターの衣装候補を容易に作成できる一方，候補が人間にどのような感性的効果を与えるかを評価する方法は十分に整理されていない．本研究では，元画像とプロフィールから形成される基準印象と衣装変更後の印象を同一参加者から取得し，印象差，意外性，キャラクターらしさ，似合い，違和感とギャップ萌えとの関係，および色・服装ジャンル・シルエット等の具体的な服装変換要因の寄与を分析する．さらに，人間評価を基準として複数の既存Vision-Language Model（VLM）に同一刺激のギャップ萌えを評価させ，人間との一致度と入力文脈の効果を比較する．30刺激を用いた探索的検証を通して，人間がどのような衣装変化にギャップ萌えを感じるかを明らかにするとともに，VLMがその感性判断をどの程度再現できるかを評価する枠組みを構築する．\\
 ＜キーワード＞ギャップ萌え，キャラクターデザイン，感性評価
\vspace{2.5em}
\end{abstract}]
\vspace{1em}

%--- 本文開始 ---
\section{はじめに}
近年，生成AIの発展により，キャラクターの衣装候補を容易に生成・編集できるようになった．衣装はキャラクターの性格や雰囲気の印象に関わる要素であり，普段とは異なる衣装によって新たな魅力が表れることがある．このように，先に形成された人物像とは異なる一面に魅力を感じる反応は，一般に「ギャップ萌え」と呼ばれる．本研究では，キャラクターの衣装変化によって生じるギャップ萌えを対象とする．
%===============================
ギャップ萌えは，衣装変更後の見た目だけでなく，そのキャラクターについて事前に形成された人物像との関係によって生じる感性反応である．同じ衣装であっても，元のキャラクター像によって自然な一致として受け取られる場合と，意外な一面として魅力につながる場合がある．そのため，衣装変更前に形成される基準の印象，変更後の印象，両者の関係を一体として扱うことが重要である．

%印象変化後の印象だけではなく，衣装変更前に形成されていたキャラクターの印象と，衣装変化によって生じた印象の変化を考える必要がある．
%===============================
そこで本研究では，衣装変更前のキャラクター画像とプロフィールから形成される印象を「基準印象」とし，衣装変更後の印象との差異に着目する．同一参加者から衣装変更前後の印象とギャップ萌え評価を取得し，人間がどのような印象変化にギャップ萌えを感じるかを分析する．さらに，色，服装ジャンル，フォーマル度，シルエット等の具体的な服装変換要因について，どの要素がギャップ萌えを高める，または弱めるものとして知覚されたかを明らかにする．

%================================

加えて，人間の評定を刺激ごとの基準値として集約し，複数の既存Vision-Language Model（VLM）に同一のギャップ萌え評価を行わせる．人間評価とVLM評価の一致度を比較するとともに，元画像やプロフィールといった基準キャラクター情報を入力へ加えたときの変化を調べることで，VLMが人間のギャップ萌え判断をどの程度再現できるかを検証する．

%また，近年では画像を理解できるVision-Language Model（VLM）の発展により画像から人間の印象や感性に関する評価を行うことも可能になっている．そこで本研究では，人間によるギャップ萌え評価をVLMによる評価を比較し，VLMが人間とどの程度同様の評価を行えるのかを調べる．さらに，衣装変化後の画像だけでなく，衣装変化前のキャラクター画像やプロフィールを与えた場合についても比較し，もともとのキャラクター情報がVLMのギャップ萌え評価に与える影響を検討する．

%===============================

本研究では次の三つの研究質問を設定する．
\begin{itemize}
 \item RQ1：衣装変更によって知覚されるキャラクター印象の変化は，ギャップ萌え評価とどのように関係するか．
 \item RQ2：どのような服装要素の変化が，ギャップ萌えを高める／弱める要因として人間に知覚されるか．
 \item RQ3：既存VLMは人間のギャップ萌え評価をどの程度再現でき，元画像やプロフィールという基準キャラクター情報はその判断にどの程度寄与するか．
 %RQ3：既存VLMによるギャップ萌え評価は，人間による評価とどの程度一致するか．また，衣装変更前のキャラクター情報はVLMの評価にどのような影響を与えるか．
\end{itemize}
本稿では，10人のオリジナルキャラクターと各3種類の衣装変化からなる30刺激を用いて，主観評価の収集から人間感性分析，服装要因分析，VLM評価までを一貫して検証する．この検証を，ギャップ萌え評価ベンチマークの構築へ向けた初期段階として位置付ける．

%本研究では，これらの問いを明らかにすることで，キャラクターの衣装変化によって生じるギャップ萌えを評価するための基礎的な知見を得るとともに，今後のギャップ萌え評価ベンチマークの構築につなげることを目指す．

%===============================

\section{関連研究}
\subsection{ギャップ萌えと衣装による人物印象の変化}
Shibuyaらは，ゲームキャラクターへの嗜好に関するインタビューから，既存の人物像とは異なる隠れた側面の発見が魅力となる「gap-moe」を報告している~\cite{1}．この知見は，ギャップ萌えが事前に形成された人物像と新たに提示された一面との関係に依存することを示唆している．
%===============================

%一方
Chengらは，仮想キャラクターの衣装について色，デザイン，種類を変化させ，衣装の違いがキャラクターから受ける性格印象に影響することを主観評価により示している~\cite{2}．本研究では，この衣装と人物印象の関係をさらにギャップ萌えへ接続し，衣装変更前後の印象差とギャップ萌えの強さを同一参加者から対応付けて取得する．また，色，服装ジャンル，フォーマル度，シルエット，装飾，露出，柄・素材感という具体的な服装変換要因について，ギャップ萌えへの知覚的な寄与を収集する．
%この研究から，衣装はキャラクターの人物印象に影響する重要な要素であることが分かる．
%しかし，衣装を変更したときに生じる人物印象の変化が，ギャップ萌えとどのように関係するのかは明らかになっていない．また，どのような衣装の変化がギャップ萌えにつながるのかについても，十分に整理されていない．そこで本研究では，衣装変化前後の人物印象とギャップ萌えの評価を取得し，両者の関係を分析する．さらに，色，服装ジャンル，フォーマル度，シルエット，装飾，露出，柄・素材感などの衣装要素にも着目する．

%===============================


\subsection{人間による主観評価を用いた画像評価}
人間の主観を画像評価の基準として利用する研究では，美的評価や生成画像品質の分野で人間評定を集約し，モデルの評価基準として用いる枠組みが検討されている．Hayashiらは，芸術作品，ファッション，風景の3ドメインを同一評定者が評価するXPASSを構築し，美的評価の個人差とドメインをまたぐ嗜好の共通性を分析している~\cite{3}．この研究は，主観評価に個人差が存在する状況でも，複数評定者の評価を用いて集団傾向と個人差の双方を扱えることを示している．

% 人間の主観を画像評価の指標として活用する研究では，美的評価や生成画像の品質評価において，人間による評定を集約し，モデル評価の基準に取り込む枠組みが検討されている．Hayashiらは，芸術作品・ファッション・風景の3ドメインを同一の評定者が評価するXPASSを整備し，美的評価における個人差と，ドメインを横断した嗜好の共通性を分析した~\cite{3}．この研究は，主観評価に個人差がある場合でも，複数の評定者の結果を用いることで，集団としての傾向と個人差の両方を扱えることを示している．

%===============================


Hirakawaらは，Virtual Try-On画像に対する大規模な人間評価からVTON-QBenchを構築し，人間の主観品質スコアとのSRCC，PLCC，pairwise accuracyを用いて評価手法を比較している~\cite{4}．また，生成画像と入力画像との関係を明示的に扱うことで，人間評価との整合性を高めている．この考え方は，元キャラクター像と衣装変更後画像の関係によって生じるギャップ萌えの評価にも対応する．

% (no)Hirakawaらは，大規模な人手評価に基づいてVirtual Try-On画像のためのVTON-QBenchを構築し，人間の主観的品質スコアに対するSRCC，PLCC，およびpairwise accuracyを指標として各評価手法を比較している~\cite{4}．さらに，生成画像と入力画像の対応関係を明示的に取り込むことで，人間評価との一致度を向上させている．この枠組みは，元のキャラクター像と衣装変更後の画像の関係から生じるギャップ萌えの評価にも適用できる．

% Hirakawaらは，Virtual Try-Onによって生成された画像を対象として大規模な人間評価を行い，その評価を用いたVTON-QBenchを構築している~\cite{4}．さらに，人間による品質評価との相関や，2つの画像のどちらが良いかを正しく判断できる割合などを用いて，画像評価手法を比較している．また，生成された衣装画像だけでなく，入力画像との関係を考慮することで，人間の評価に近い評価が可能になることを検討している．



%===============================


本研究は，ギャップ萌えを基準キャラクターと衣装変化の関係に依存する感性反応として扱い，第一に人間評価からその成立傾向と服装要因を分析し，第二にその人間評価を基準として複数VLMの感性判断を比較する．これにより，衣装変化に対する人間感性の分析とVLM評価を同一の評価データ上で接続する．

% これらの研究では，人間による主観評価を用いて画像評価手法を検証している．本研究でも，人間によるギャップ萌え評価を用いてVLMの評価を検証する．特に，衣装変更後の画像だけを見た場合と，衣装変更前のキャラクター画像やプロフィールを含めた場合を比較し，キャラクターのもともとの情報がVLMのギャップ萌え評価に与える影響を調べる．

\section{提案手法}
\subsection{全体構成}
本研究の枠組みは，(1)キャラクターの基準印象の取得，(2)衣装変化後の印象・心理反応・ギャップ萌えの取得，(3)人間評価の集約と感性分析，(4)服装変換要因の分析，(5)複数VLMによる同一刺激の評価とHuman--VLM Alignmentの比較，の五段階から構成される．人間評価は刺激ごとの基準値として保持し，VLMには元画像，プロフィール，衣装変化画像を入力する．これにより，人間の感性構造を分析した上で，同じ刺激に対するVLMの判断との一致度を評価する．

%本研究では，キャラクターの衣装を変更したときに生じる印象の変化と，その変化に対するギャップもrの評価との関係を調べる．また，衣装のどのような変化がギャップ萌えに関係するのかを分析し人間による評価とVision-Language Model（VLM）による評価を比較する．

%本研究の流れを図○に示す．まず衣装を変更する前のキャラクター画像とプロフィールから，キャラクターに対する印象を参加者に評価してもらう．次に衣装を変更した画像を提示し，変更後の印象とギャップ萌えの程度を評価してもらう．得られた評価から衣装変更による印象の変化とギャップ萌えとの関係，およびギャップ萌えに関係する衣装要素を分析する．最後に人間が評価したものと同じ衣装変化についてVLMにも評価させ，人間による評価との一致を調べる．

%===============================

\subsection{主観評価の収集}
参加者には研究用に作成したオリジナルキャラクターを提示する．各キャラクターは，全身の元画像と，年齢，性格，趣味等を簡潔に記したプロフィールから構成する．まず衣装変化を見る前に，I1「かわいい」，I2「かっこいい」，I3「落ち着いている」，I4「活発である」，I5「大人っぽい」，I6「親しみやすい」の6項目を7件法で回答してもらい，参加者ごとの基準印象を取得する．

次に，同じキャラクターの元画像と衣装変更後画像を提示し，再度I1--I6を回答してもらう．加えて，G1「元の印象と異なる」，G2「意外だと感じる」，G3「同じキャラクターらしさが保たれている」，G4「似合っている」，G5「魅力が増した」，G6「違和感を感じる」を7件法で取得する．ギャップ萌えについては，GM1「この衣装変化にギャップ萌えを感じた」を主要評価項目とし，GM2「元の印象との違いそのものに魅力を感じた」，GM3「このキャラクターの意外な一面に惹かれた」を補助項目として取得する．GM1はギャップ萌えそのものを直接問う基準値として用い，GM2とGM3はその構成概念を補助的に確認するために用いる．

衣装要因に関する質問は，I，G，GMの回答を確定した後に提示する．これにより，色やシルエット等の観点を先に示すことによる全体評価へのプライミングを抑える．服装要因は，F1 色・色調，F2 服装ジャンル・テイスト，F3 フォーマル度，F4 シルエット・形状，F5 装飾・小物，F6 肌の露出・身体の見せ方，F7 柄・素材感の7領域とする．参加者は元衣装から「変化したと感じた要因」を複数選択し，選択した要因について，その変化がギャップ萌えをどの方向にどの程度影響したと感じたかを$-2$から$+2$で回答する．


% 参加者には，研究用に作成したオリジナルキャラクターを提示する．各キャラクターについて，全身の元画像と，年齢，性格，趣味などを簡潔に記したプロフィールを用意する．

% まず，衣装変更後の画像を見る前に，キャラクターの元画像とプロフィールを提示する．参加者には，キャラクターから受ける印象について，I1「かわいい」，I2「かっこいい」，I3「落ち着いている」，I4「活発である」，I5「大人っぽい」，I6「親しみやすい」の6項目を7件法で評価してもらう．ここで得られた評価を，衣装変更前のキャラクターに対する印象として扱う．

% 次に，同じキャラクターについて，元画像と衣装変更後の画像を提示する．参加者には，再びI1--I6を評価してもらい，衣装変更後にキャラクターの印象がどのように変化したかを調べる．さらに，G1「元の印象と異なる」，G2「意外だと感じる」，G3「同じキャラクターらしさが保たれている」，G4「似合っている」，G5「魅力が増した」，G6「違和感を感じる」の6項目を7件法で評価してもらう．

% ギャップ萌えについては，GM1「この衣装変化にギャップ萌えを感じた」を主要な評価項目とする．また，GM2「元の印象との違いそのものに魅力を感じた」，GM3「このキャラクターの意外な一面に惹かれた」の2項目も補助的に評価してもらう．GM1はギャップ萌えそのものを直接評価する項目として用い，GM2とGM3は，ギャップ萌えに含まれる「元の印象との違い」や「意外な一面への魅力」との関係を確認するために用いる．

% 最後に，衣装の変化について評価してもらう．この質問は，I，G，GMの評価を終えた後に提示することで，衣装の具体的な要素に注目することが先の評価に影響しないようにする．

% 衣装の変化は，F1「色・色調」，F2「服装ジャンル・テイスト」，F3「フォーマル度」，F4「シルエット・形状」，F5「装飾・小物」，F6「肌の露出・身体の見せ方」，F7「柄・素材感」の7項目に分ける．参加者には，元の衣装と比べて「変化したと感じた要因」を複数選択してもらう．その後，選択した各要因について，その変化がギャップ萌えにどの程度影響したと感じたかを$-2$から$+2$で評価してもらう．


%===============================
\subsection{人間基準スコアと評価信頼性}
刺激$s$に対する参加者$u$のGM1評定を$g_{s,u}$とする．同じ刺激を評価した参加者集合を$U_s$とし，人間集団による刺激レベルのギャップ萌え基準値$H_s$を
\begin{equation}
 H_s=\frac{1}{|U_s|}\sum_{u\in U_s}g_{s,u}
\end{equation}
と定義する．$H_s$は，本実験条件において複数参加者から得られた集団的な主観評価の代表値として扱う．平均値に加えて，標準偏差，中央値，95\%信頼区間，各評定値の分布を保存し，刺激ごとの評価のばらつきも報告する．

人間基準値の再現性を確認するため，参加者を各刺激の割当を保ったまま二群に分け，各群から独立に刺激平均$H_s^{A}$，$H_s^{B}$を求める．この分割を繰り返し，Spearman順位相関係数（SRCC），Pearson相関係数（PLCC），平均絶対誤差（MAE）の分布を求める．このHuman--Human Alignmentにより，個人差を含む主観評価において，刺激間の相対的な評価傾向がどの程度再現されるかを確認する．

\subsection{印象変化とギャップ萌えの分析}
衣装変更前の6次元基準印象を$\mathbf{b}_{c,u}$，変更後印象を$\mathbf{t}_{c,j,u}$とする．参加者$u$がキャラクター$c$の衣装$j$に対して知覚した印象変化量$d_{c,j,u}$を，各軸の差のユークリッド距離を理論上の最大距離で正規化して
\begin{equation}
 d_{c,j,u}=\frac{\sqrt{\sum_{m=1}^{6}(t_{c,j,u,m}-b_{c,u,m})^2}}{6\sqrt{6}}
\end{equation}
と定義する．$0\le d_{c,j,u}\le1$であり，値が大きいほど参加者本人が知覚した人物印象の変化が大きいことを表す．

RQ1では，まず$d$とGM1の線形関係を確認し，次に$d^2$を含むモデルとの比較から，印象変化量とギャップ萌えの関係が単調であるか，中程度の変化で高くなる非線形関係を持つかを検討する．さらにG2，G3，G4，G6を含め，意外性，キャラクターらしさ，似合い，違和感がGM1とどのように関連するかを分析する．とくに$d\times G3$の交互作用を用いて，「元の印象とは異なるが同じキャラクターとして受け入れられる」という状態とギャップ萌えとの関係を調べる．

GM1は1--7の順序尺度であり，同一参加者が複数刺激を評価するため，主要分析には累積ロジット型の階層順序回帰モデルを用いる．参加者，キャラクター，刺激による変動を階層効果として扱い，参加者ごとの評定傾向と刺激に対する反応を分離する．また，$d$とG1，GM1とGM2・GM3の関連を補助的に確認し，計算された印象距離と主観的な印象差，主要GM評価と補助GM評価の対応を確認する．

\subsection{服装変換要因の分析}
要因$k$が変化として選択されたかを$z_{s,u,k}\in\{0,1\}$，選択時の寄与評定を$e_{s,u,k}\in\{-2,-1,0,+1,+2\}$とする．各要因について，変化として認識された頻度を表す選択率$\pi_k$，選択された場合にどちらの方向へ寄与したと知覚されたかを表す平均寄与$\mu_k$を
\begin{equation}
 \pi_k=\frac{\sum_{s,u}z_{s,u,k}}{\sum_{s,u}1},\quad
 \mu_k=\frac{\sum_{s,u}z_{s,u,k}e_{s,u,k}}{\sum_{s,u}z_{s,u,k}}
\end{equation}
として求める．さらに，要因が選択されなかった回答を0として扱う平均符号付き寄与を算出する．これにより，「よく変化として認識される要因」と「認識されたときにギャップ萌えを高める／弱める方向に関連した要因」を分離して記述する．信頼区間は参加者単位のクラスタブートストラップによって求める．一つの衣装内で複数属性が同時に変化するため，本結果は参加者がギャップ萌えに関係したと知覚した要因として解釈する．属性ごとの因果効果は，今後，個別属性を独立に操作する実験によって検証できる．

\subsection{VLMによるギャップ萌え評価}
RQ3では，複数の既存VLMに対し，人間と同じ刺激についてギャップ萌えの強さを評価させる．VLMにはキャラクター情報と衣装変更後画像を与え，人間の評定結果は比較用の基準値として保持する．主評価タスクでは，人間アンケートと同じギャップ萌えの定義を提示し，「この衣装変化にギャップ萌えを感じる」程度を1--7の整数で出力させる．モデル間の公平性を保つため，プロンプト，画像順，出力形式，生成設定を統一し，構造化形式でGM1のみを回答させる．G1--G6の評価は主評価とは独立した診断タスクとして実行し，GM1評価へのプライミングを避ける．

ギャップ萌えが基準人物像との関係に依存するかを調べるため，次の三条件を比較する．
\begin{itemize}
 \item C0：衣装変更後画像のみ．
 \item C1：元画像＋衣装変更後画像．
 \item C2：元画像＋プロフィール＋衣装変更後画像．
\end{itemize}
C0からC1への改善は元キャラクターの視覚情報を参照する効果，C1からC2への改善は画像だけでは表現されない人物設定を加える効果として解釈する．主条件はC2とし，C0とC1は入力文脈の寄与を調べるablationとして扱う．

\section{検証実験}
\subsection{刺激とアンケート条件}
参加者が既存作品に関する知識を持たない条件で評価できるよう，研究用のオリジナルキャラクター10人を用いる．各キャラクターは全身の立ち絵と短いプロフィールによって提示する．各キャラクターに対して3種類の衣装変化画像を用意し，計30組を評価対象とする．衣装以外の影響を抑えるため，顔，髪型，髪色，体格，ポーズ，構図，背景および画風は可能な限り保持し，主に衣服と関連アクセサリを変更する．

3衣装は，(1)元の人物像との整合性を重視したBaseline-consistent，(2)元の人物像と異なる要素を導入しつつ同一人物性を維持することを狙ったIntended Gap-Moe，(3)別方向またはより強い印象変化を与えるStrong/Alternative Contrastとして設計する．これらは刺激作成時の設計条件であり，分析上の正解ラベルには用いない．どの衣装が実際にギャップ萌えを生むかは人間評定から決定する．

Webアンケートでは，各参加者が全10キャラクターについて衣装変更前の基準印象を回答した後，各キャラクターについて3衣装のうち1衣装ずつ，計10刺激を評価する．参加者を3群に分け，各キャラクターの3衣装が参加者全体で可能な限り均等に評価されるよう割り当てる．提示順はランダム化する．

\subsection{回答品質とHuman Ground Truth}
解析対象は，調査を完了し，対象版のアンケートに回答し，割当と保存データの整合が確認できた回答とする．不完全回答，重複，データ破損等は解析対象から除外する．極端に短い回答時間や同一値の連続回答等は品質フラグとして保持し，フラグ対象の有無による感度分析を行う．この処理により，回答品質を確認しながらギャップ萌えの個人差を分析対象として保持する．

有効回答確定後，各刺激についてHuman Gap-Moe Score $H_s$を算出し，Human--Human Alignmentを求める．これにより，VLM比較に先立って，人間集団内でどの程度安定した刺激順位が得られるかを確認する．

\subsection{人間感性と服装変換要因の評価}
RQ1について，$d$の一次項のみのモデルと$d^2$を含むモデルを比較し，印象変化量とGM1の関係形状を評価する．続いてG2，G3，G4，G6および$d\times G3$を含む階層順序回帰を用い，各感性要因との条件付き関連を評価する．参加者間の評定傾向を考慮し，同一参加者内で普段より高く／低く感じた変化とGM1との対応を重視して解釈する．

RQ2について，各F要因の選択率，選択時平均寄与，平均符号付き寄与と95\%信頼区間を算出する．これらを並べることで，色やシルエット等が「変化として認識されやすいか」と「ギャップ萌えを高める／弱める要因として知覚されやすいか」を区別する．さらに，基準印象I1--I6と各F要因の符号付き寄与との関係を探索的に可視化し，どのような人物像と衣装変化の組合せで反応が異なるかを検討する．

\subsection{VLM評価プロトコル}
評価対象には，画像と言語を同時に入力できる複数の既存VLMを用いる．使用モデル名とバージョンは実験時点で固定し，すべてのモデルに同一の日本語プロンプト，同一の画像，同一の出力スキーマを使用する．温度は0または利用可能な最小値とし，各刺激は独立したセッションとして評価する．出力形式はGM1の整数値のみを含むJSON形式とし，形式不正時のみ内容を変更しない再出力を1回許可する．

主条件C2について全30刺激を評価した後，C0，C1を同一設定で実行する．また，診断分析用としてC2条件でG1--G6を別セッションから取得する．生の応答，構造化後のスコア，モデル名，バージョン，実行設定を保存し，再現可能な評価ログを残す．

\subsection{Human--VLM Alignmentの評価指標}
主評価指標は，人間平均$H_s$とVLM予測$\hat{H}_{m,s}$のSRCCとする．SRCCはスコアの絶対的な尺度利用の違いに比較的頑健であり，人間がギャップ萌えを高く評価した刺激をVLMも高く順位付けできるかを表す．補助指標としてPLCCとMAEを報告する．

さらに，同一キャラクター3衣装の相対順位を評価する．人間平均で高い衣装をVLMも高く評価した割合をpairwise ranking accuracyとして求め，各キャラクターで最も高い$H_s$を持つ衣装をVLMも1位に選択できた割合をTop-1 selection accuracyとして報告する．信頼区間はキャラクター単位で再標本化するブートストラップによって求め，同一キャラクター内の3衣装を独立な標本として扱わない．

\subsection{VLMの誤差分析}
VLMが人間判断と異なる理由を具体化するため，主条件C2におけるG1--G6についても人間平均とVLM出力のSRCCおよびMAEを算出する．これにより，例えば意外性の判断は一致する一方，キャラクターらしさや違和感の判断が一致しないといったモデルごとの傾向を分析する．

また，$|\hat{H}_{m,s}-H_s|$が大きい刺激を抽出し，人間のG1--G6およびF1--F7の集計結果と照合する．人間が「意外だがキャラクターらしくない」と評価した刺激をVLMが高く評価する場合など，総合的なGM1の誤差を具体的な感性要因の違いとして記述する．この分析により，モデル間の総合的な一致度に加えて，人間のギャップ萌え判断のどの側面に一致または不一致が生じるかを検討する．

\subsection{結果}
% TODO: アンケート収集・VLM評価完了後に，以下を挿入する．
% 1) 有効参加者数，回答数，各刺激の評定数
% 2) Human--Human Alignment（split-half SRCC, PLCC, MAE）
% 3) RQ1：d, d^2, G2, G3, G4, G6, d×G3 の結果
% 4) RQ2：F1--F7 の選択率・平均寄与・平均符号付き寄与
% 5) RQ3：各VLMのC2におけるSRCC, PLCC, MAE, Pairwise, Top-1
% 6) C0/C1/C2 ablation
% 7) G1--G6診断および代表的失敗例
\textbf{本節には，アンケート収集およびVLM評価完了後に定量結果を挿入する．}

\subsection{考察}
本研究の第一の目的は，衣装変化に対するギャップ萌えを，元の人物像との関係に基づく感性反応として捉えることである．RQ1の結果から，印象変化量そのものの効果に加えて，意外性，キャラクターらしさ，似合い，違和感との関連を確認することで，人間がどのような状態をギャップ萌えとして受け止めるかを整理できる．とくに印象変化量とキャラクターらしさの交互作用が確認される場合，「異なる一面であること」と「同じ人物として受け入れられること」の両立が重要であるという解釈につながる．

第二に，F1--F7の要因分析は，ギャップ萌えという抽象的な感性反応を，衣装設計で操作可能な要素へ対応付ける役割を持つ．高いGM1を得た刺激について，どの要因が変化として認識され，正の寄与として回答されたかを整理することで，色，ジャンル，フォーマル度，シルエット等のどの変化が人間の知覚上関係したかを説明できる．複数属性が同時に変化する本実験からは知覚された関連を整理し，属性ごとの因果効果は今後の独立操作実験へ接続する．

第三に，Human--VLM Alignmentは，VLMがギャップ萌えという文脈依存の感性判断をどの程度人間に近く再現できるかを評価する．C0，C1，C2の比較によって元画像やプロフィールを追加した際に人間との一致度が向上する場合，変更後画像単体の外見だけではなく，元のキャラクター像を参照することがギャップ萌え判断に有効であることを示せる．差が小さい場合には，今回の刺激構成において変更後画像から得られる情報が主要な判断材料となっている可能性を示す．この比較により，ギャップ萌え評価における基準文脈の役割を共通プロトコル上で定量化する．

本検証では10キャラクター30刺激を用いて，人間評価の収集，人間感性の特徴付け，VLMとの比較までを一貫して実施し，ギャップ萌え評価ベンチマーク構築に向けた初期プロトコルを検証する．今後はキャラクター数と衣装変化の種類を拡張し，評価の安定性と多様性を高めたベンチマークへ発展させる．さらに，参加者属性や過去評価を取り入れ，個人ごとのギャップ萌え感性を扱う評価へ拡張する．

\section{まとめと今後の課題}
本研究では，キャラクターの元画像とプロフィールから形成される基準印象と，衣装変更後に知覚される印象・心理反応・ギャップ萌えを同一参加者から取得し，人間がどのような衣装変化にギャップ萌えを感じるかを分析する枠組みを示した．また，色，服装ジャンル，フォーマル度，シルエット等の具体的な服装変換要因について，変化として知覚される頻度とギャップ萌えへの正負の寄与を分離して分析する．さらに，刺激ごとの人間平均を基準値として複数の既存VLMに同一評価を行わせ，Human--VLM Alignmentと入力文脈のablationから，人間の感性判断をVLMがどの程度再現できるかを評価する．30刺激による本検証を，ギャップ萌え評価ベンチマーク構築に向けた第一段階とする．今後は刺激数とキャラクター多様性を拡張し，人間評価の安定性，VLM間比較の再現性，個人差を含む感性評価をより大規模に検証する．

\begin{thebibliography}{9}
  \footnotesize

  \bibitem {1} A. Shibuya, H. Okura, A. Shoun, and N. Asou: ``Male and Female Game Players' Preferences for Game Characters and Real-world Personalities in Japan,'' Proceedings of DiGRA 2019 Conference: Game, Play and the Emerging Ludo-Mix, 2019.
  \bibitem {2} Y. Cheng and Y. Wang: ``Evaluating the Effect of Outfit on Personality Perception in Virtual Characters,'' Virtual Worlds, Vol.3, No.1, pp.21--39, 2024.
  \bibitem {3} T. Hayashi, H. Takahara, C. O. Mawalim, H. Narimatsu, A. Kimura, S. Kumano, and S. Okada: ``XPASS-Vis: A Dataset for Cross-Domain Personalized Image Aesthetic Assessment,'' arXiv preprint arXiv:2606.15629, 2026.
  \bibitem {4} Y. Hirakawa, T. Wada, R. Shimizu, T. Furusawa, Y. Saito, R. Araki, T. Chen, F. Mo, and Y. Aoki: ``Reference-Free Quality Assessment for Virtual Try-On via Human Feedback,'' arXiv preprint arXiv:2603.13057, 2026.

\end{thebibliography}
\footnotesize

\noindent{杉原 明音：}
明星大学情報学部情報学科所属．キャラクターデザインと感性評価に関する研究に取り組む．世界最高の厨二病．
\vskip\baselineskip

\noindent{山口 熙己：}
明星大学情報学部情報学科所属．画像認識，Vision-Language Model，人間評価を用いたAIシステムの研究に取り組む．
\vskip\baselineskip

\noindent{植木 一也：}
明星大学情報学部情報学科所属．画像認識・機械学習に関する研究に取り組む．

\end{document}
